% !TEX root =  Lecture11.tex


%%%%%%%%%%%%%%%%%%%%%%%%%%%%%%%%%%%%
% Recap/Agenda
%%%%%%%%%%%%%%%%%%%%%%%%%%%%%%%%%%%%
\begin{frame}
\frametitle{Module 2: Statistical Inference}
    \begin{itemize}
        \item \hl{Previously:} inference from \emph{formulas}: the CLT, $z$ and $t$, two-sample and paired data. All of it about means and proportions.
        \item \hl{This time:} what if both variables are \emph{categorical}?
          \begin{itemize}
              \item two-way tables, and what independence predicts: \emph{expected counts};
              \item the $X^2$ statistic, one number for a whole table;
              \item its null distribution, the same two ways: shuffle the labels, or use the $\chi^2$ model;
              \item conditions, and reading the residuals rather than just the p-value.
          \end{itemize}
        \item \hl{Reading:} IMS Chapter 18.
        \item \hl{Homework for today:}
        \item \hl{Homework for next class:}
    \end{itemize}
\end{frame}


%%%%%%%%%%%%%%%%%%%%%%%%%%%%%%%%%%%%
% Learning objectives:
%%%%%%%%%%%%%%%%%%%%%%%%%%%%%%%%%%%%
\begin{frame}
    \frametitle{Lecture Objectives}
    \goals{%
    \begin{itemize}
    \item state what independence means for a two-way table, and turn it into expected counts;
    \item compute $X^2$ and say what each piece of it is doing;
    \item obtain its null distribution by shuffling \emph{and} from the $\chi^2_{df}$ model, and explain why they agree;
    \item check the conditions, and know which road survives when they fail;
    \item interpret the result: the p-value for \emph{whether}, the standardized residuals for \emph{where}.
    \end{itemize}}
    \objectives{%
    \begin{itemize}
    \item \textbf{(M2, LO4)} Testing Independence
    \end{itemize}}
\end{frame}
